% Ad Atticum 10.17 --- headnote
% ad-atticum-10-17, 16 May 49 BC, Cumae

\textit{Cicero to Atticus, written from the Cuman villa
on the seventeenth day before the Kalends of June 49 BC ---
16 May (the manuscript dateline: \textit{Scr.\ in Cumano
xvii K.\ Iun.\ a.\ 705 (49)}). The letter is a short
operational dispatch in four brief sections, written
while the household is still preparing the crossing
that has been hanging fire since the start of the month.
Hortensius the younger, governor of the neighbouring
coast, has called on Cicero on the day before the Ides
(14 May) --- with a letter to him already written, an
awkward fact Cicero notes in passing --- and is being
extravagantly obliging, an \textit{[Greek: ekten\=eian]},
``devotion,'' that Cicero plainly means to spend.
Atticus's freedman Serapion has come in close behind
with a letter that confirms what Cicero already knew
about him from an earlier note: learned and honest,
and worth keeping. Cicero now means to take the ship
he sails on, and Serapion himself as fellow-passenger.}

\textit{The rest is brief housekeeping. The eye trouble
(\textit{lippitudo}) keeps flaring up just enough to
slow Cicero's writing; he is glad to hear Atticus is
finally clear of his old illness and the new bouts that
followed it. He wishes Ocella were here, and notes that
the chief obstacle now is the equinox, which has run
rough; if it settles, he hopes Hortensius continues as
generous as he has been. Section 4 answers a small note
of pique from Atticus: Cicero had assumed that Atticus
held an imperial \textit{diploma} (an official
travel-warrant of the kind a magistrate issued to
authorise the use of the cursus publicus), partly
because Atticus had spoken of travel and partly because
he had drawn one out for his slaves; Cicero says so and
asks for any further news, however small. The letter
closes with a dateline of its own: xvii K.\ Iun., 16
May. Two daggered cruxes in the manuscript (one in \S1,
one in \S3) are preserved in this translation as \dag\
markers; the sense given is the most natural reading
of the corrupted text.}
